\subsection{\label{CayleyHamilton}Cayley-Hamilton method and some results}

Cayley-Hamilton method is widely used in the calculation of smearing, here is some results.

\subsubsection{\label{CayleyHamiltonMethod}Cayley-Hamilton method}

Cayley-Hamilton theory demonstrated, for square matrices,
\begin{equation}
\begin{split}
&p(M)=0,
\end{split}
\end{equation}
where, $p(x)=0$ is the characteristic equation, i.e., $p(x)=\det(M-x I)$ (`p' means polynomial).

It is used to calculate the function of a $N\times N$ matrix, 
\begin{equation}
\begin{split}
&f(M)=\sum _{i=0}^{N-1} f_i M^i,
\end{split}
\end{equation}

To determine the coefficients, note that, when $\lambda _j$ is eigen-value of $M$, $f(\lambda _j)$ is eigen-value of $f(M)$. So, we can,

(1) Calculate $f(\lambda _j)$ use $f$.

(2) Calculate $f(\lambda _j)$ use $f(M)=\sum _{i=0}^{N-1} f_i M^i$ with TBD coefficients, $f_i$.

(3) Solve the equation $f(\lambda _j)=\sum _{i=0}^{N-1} f_i \lambda _j^i$ to determine $f_i$.


\subsubsection{\label{CayleyHamilton3by3}3 x 3 matrix}

\begin{equation}
\begin{split}
&f(Q)=f_0+f_1Q+f_2Q^2,
\end{split}
\end{equation}
with,
\begin{equation}
\begin{split}
&f(\lambda_0)=f_0+f_1\lambda_0+f_2\lambda_0^2,\\
&f(\lambda_1)=f_0+f_1\lambda_1+f_2\lambda_1^2,\\
&f(\lambda_2)=f_0+f_1\lambda_2+f_2\lambda_2^2,\\
\end{split}
\end{equation}
where $\lambda_{0,1,2}$ are eigen-values of $Q$.


\subsubsection{\label{CayleyHamiltonExponetial}Exponential}


\subsubsection{\label{CayleyHamiltonSqrt}Sqrt}